Understanding the brain at the level of individual synaptic connections requires
mapping the complete wiring diagram --- the \emph{connectome} --- of a neural
circuit. Modern connectomics workflows acquire volumetric image stacks of brain
tissue using electron microscopy (EM) or X-ray nano-holotomography (XNH), then
apply automated deep-learning segmenters to delineate individual neurons within
those volumes~\cite{januszewski2018high,lin2021pytorch}. This process has enabled
landmark reconstructions at increasingly large scales, culminating in full-cortex
and whole-brain datasets containing tens of thousands of neurons and hundreds of
millions of synapses~\cite{microns2025}.

Despite rapid progress in segmentation quality, automated segmenters remain
imperfect. They produce two types of errors: \emph{merge errors}, where
two distinct neurons are incorrectly fused into one segment, and \emph{split errors},
where a single neuron is fragmented into multiple disconnected pieces. Split
errors are particularly prevalent in elongated structures such as axons, where
slight imaging artifacts, staining variation, or object boundaries that are
difficult to resolve cause segmenters to break a continuous process into two or
more isolated fragments. At the scale of a full cortex reconstruction, a single
well-segmented volume may contain hundreds of thousands of split errors that must
be identified and corrected before the connectivity graph is reliable
enough for downstream neuroscience~\cite{plaza2018analyzing}.

Manual proofreading --- having human experts identify and repair segmentation
errors --- is the bottleneck in every large-scale connectomics project. A skilled
proofreader can correct perhaps a few hundred merge decisions per hour; a modern
dataset may require millions. Automated proofreading tools that propose high-recall
merge corrections, leaving only ambiguous cases for human review, are therefore
essential infrastructure for the field~\cite{dorkenwald2022cave}.

\section*{The Problem This Work Addresses}

This project targets \textbf{split detection and correction proposal} in
post-segmentation connectomics pipelines. Given a voxel-wise segmentation
volume (the output of an automated segmenter), the pipeline must:

\begin{enumerate}
  \item Identify all pairs of fragments that are likely pieces of the same neuron
        (\emph{split candidates}).
  \item Score and filter these pairs conservatively, accepting only those for which
        evidence strongly supports a merge.
  \item Explicitly flag uncertain pairs rather than forcing them to a binary decision.
  \item Produce outputs in formats consumable by downstream proofreading tools.
\end{enumerate}

The domain of primary interest is \textbf{myelinated axons in mouse white matter},
as captured by XNH imaging --- the XPRESS challenge dataset~\cite{xpress2023}.
Myelinated axons are predominantly tubular structures with consistent cross-sectional
geometry, making geometric alignment and curvature features informative proxies
for merge plausibility. This is in contrast to dendritic arbors, where branching
angles can be nearly orthogonal and geometric features are less discriminative.

\section*{Contributions}

The main contributions of this work are:

\begin{enumerate}
  \item \textbf{A complete, modular post-segmentation pipeline} implemented in
        Python, with YAML-configurable parameters, six processing stages, and
        multi-format export. The pipeline is fully open-source and validated by
        428 automated unit and integration tests.

  \item \textbf{The skeleton-node graph construction method}, which indexes every
        node from every TEASAR skeleton in a single KD-tree and queries the full
        node set for proximity-based edge creation. This exposes interior splits
        along long axons that an endpoint-only graph cannot represent, producing
        a 47-percentage-point improvement in oracle coverage on XPRESS (from 23\%
        to 79.8\% in a single change).

  \item \textbf{Distance-conditioned composite scoring}, which switches the
        candidate score weighting from a proximity-dominated formula to a
        proximity-free, alignment-continuity-weighted formula for long-range
        pairs (gap $> 1000$~nm). This prevents the suppression of genuine long-range
        splits whose proximity score decays to near zero despite strong directional
        alignment.

  \item \textbf{A three-outcome validation model} (ACCEPT / REJECT / AMBIGUOUS)
        with seven composable, configurable rules. The AMBIGUOUS outcome
        explicitly preserves uncertainty rather than forcing binary decisions,
        making the pipeline suitable for integration with human-in-the-loop
        proofreading workflows.

  \item \textbf{Quantitative evaluation on two benchmarks}: CREMI (F1 = 0.952,
        Precision = 1.000) and XPRESS (Recall = 0.994 on training, Recall = 0.885
        on a fully held-out validation volume), with 23 documented experiments
        tracing the development of each improvement.
\end{enumerate}

\section*{Report Organization}

Chapter~\ref{ch:background} provides background on connectomics imaging, automated
segmentation, and the XPRESS challenge. Chapter~\ref{ch:methodology} describes
the pipeline design and all algorithmic components in detail.
Chapter~\ref{ch:results} presents quantitative results across all experiments.
Chapter~\ref{ch:relatedwork} surveys related proofreading and graph-construction
approaches. Chapter~\ref{ch:conclusion} summarizes the main findings and
challenges, and Chapter~\ref{ch:futurework} outlines the plan for the remainder
of the semester.
